\documentclass[11pt,a4paper]{article}

\usepackage{kotex}
\usepackage{fontspec}

% 현대적이고 부드러운 Pretendard 폰트 설정
\setmainfont{Pretendard}[
  UprightFont=* Regular,
  BoldFont=* SemiBold,
]
\setmainhangulfont{Pretendard}[
  UprightFont=* Regular,
  BoldFont=* SemiBold,
]
\setsansfont{Pretendard}[
  UprightFont=* Regular,
  BoldFont=* SemiBold,
]
\setsanshangulfont{Pretendard}[
  UprightFont=* Regular,
  BoldFont=* SemiBold,
]

\usepackage[top=22mm, bottom=22mm, left=22mm, right=22mm, headheight=22pt]{geometry}
\usepackage{amsmath, amssymb, amsfonts, mathtools}
\usepackage{xcolor}
\usepackage{enumitem}
\usepackage{titlesec}
\usepackage{fancyhdr}
\usepackage{setspace}
\usepackage[colorlinks=true, linkcolor=blue!70!black, urlcolor=blue!70!black, citecolor=blue!70!black]{hyperref}
\usepackage{bm}

% 줄간격 및 문단 설정
\setstretch{1.32}
\setlength{\parskip}{0.4em}
\setlength{\parindent}{0pt}

% 색상 정의
\definecolor{primary}{RGB}{28, 80, 135}
\definecolor{secondary}{RGB}{45, 130, 195}
\definecolor{darkgray}{RGB}{70, 70, 70}
\definecolor{rulecolor}{RGB}{210, 220, 235}

% 페이지 헤더 및 푸터 설정
\pagestyle{fancy}
\fancyhf{}
\lhead{\footnotesize\textcolor{darkgray}{\textbf{QML 연구 계획 보고서}}}
\rhead{\footnotesize\textcolor{gray}{Quantum Machine Learning Research Proposal}}
\cfoot{\small\thepage}
\renewcommand{\headrulewidth}{0.6pt}
\renewcommand{\headrule}{\hbox to\headwidth{\color{rulecolor}\leaders\hrule height \headrulewidth\hfill}}

% 섹션 넘버링 및 서식 정의
\titleformat{\section}
  {\Large\bfseries\color{primary}}
  {}{0em}{}
  [\vspace{0.2em}{\color{secondary}\hrule height 1pt}]

\titleformat{\subsection}
  {\large\bfseries\color{primary}}
  {}{0em}{}

\titleformat{\subsubsection}
  {\normalsize\bfseries\color{secondary}}
  {}{0em}{}

\titlespacing*{\section}{0pt}{1.4ex plus 0.4ex minus .2ex}{0.7ex plus .2ex}
\titlespacing*{\subsection}{0pt}{1.1ex plus 0.3ex minus .2ex}{0.4ex plus .2ex}
\titlespacing*{\subsubsection}{0pt}{0.9ex plus 0.2ex minus .2ex}{0.3ex plus .2ex}

% 리스트 여백 조정
\setlist[itemize]{leftmargin=1.5em, itemsep=0.25em, topsep=0.15em, parsep=0.1em}
\setlist[enumerate]{leftmargin=1.5em, itemsep=0.25em, topsep=0.15em, parsep=0.1em}

\begin{document}

% 문서 제목 블록
\begin{center}
    {\LARGE \textbf{\color{primary} QML 연구 계획 보고서}} \\[0.35em]
    \vspace{0.2em}
    {\color{primary}\hrule height 1.8pt}
    \vspace{0.5em}
\end{center}

\section*{I. 연구 배경 및 참고 문헌 (Background \& References)}
\addcontentsline{toc}{section}{I. 연구 배경 및 참고 문헌 (Background \& References)}

최근 양자 기계학습(QML)과 양자 시뮬레이션 분야는 노이즈가 있는 중간 규모 양자 하드웨어(NISQ)의 물리적 한계와, 최적화 과정에서 발생하는 바렌 플래토(Barren Plateaus) 현상이라는 두 가지 큰 장벽에 직면해 있습니다. 본 연구는 아래 네 편의 핵심 선행 연구가 지닌 한계점들을 분석 및 보완하여, 고전 컴퓨터의 시뮬레이션 한계를 넘어서는 2D 스핀 유리(Spin Glass) 모델 동역학 시뮬레이션 및 양자 우위 증명 패러다임을 제안합니다.

\subsection*{1. 핵심 참고 문헌 및 이론적 배경}

\begin{itemize}
    \item \textbf{Train on Classical, Deploy on Quantum (TCDQ)}
    \begin{itemize}[label=--]
        \item \textbf{내용:} 양자 생성 모델의 막대한 최적화 비용을 줄이기 위해, 고전 컴퓨터에서 손실 함수(MMD)와 그래디언트를 효율적으로 계산하여 파라미터를 훈련한 뒤, 양자 하드웨어에서는 샘플링만 수행하는 하이브리드 전략을 제시했습니다.
        \item \textbf{한계 및 극복:} 해당 연구는 교환 가능한(Commuting) 대각 게이트 위주의 IQP 회로와 특정 편향성을 가진 데이터에 국한되어 있습니다. 본 연구는 이를 \textbf{비가환적(Non-commuting) 상호작용이 포함된 범용 동역학 모델로 확장 적용}합니다.
    \end{itemize}

    \item \textbf{Backpropagating Pauli Propagation (BP-PPS)}
    \begin{itemize}[label=--]
        \item \textbf{내용:} 양자 회로의 가역성을 이용하여 중간 상태를 저장하지 않고 파울리 전파를 역전파하는 알고리즘입니다. 이를 통해 국소 관측가능량(Local observables)에 대한 해석적 그래디언트를 메모리 효율적으로 빠르게 계산할 수 있습니다.
        \item \textbf{적용:} 이 기술을 활용하여, 매우 깊은 트로터(Trotter) 진화 회로를 깊이가 얕은 하드웨어 효율적 안자츠(HEA)로 압축하는 고전적 최적화를 수행합니다.
    \end{itemize}

    \item \textbf{Limits of Quantum Generative Models with Classical Sampling Hardness}
    \begin{itemize}[label=--]
        \item \textbf{내용:} 고전적으로 샘플링이 어려운 양자 모델(Anticoncentration 분포)은 MMD와 같은 전역적인 손실 함수를 사용할 경우 평균값과 분산이 지수적으로 소실되어 훈련이 불가능해진다는 수학적 한계를 증명했습니다.
        \item \textbf{돌파구:} 본 연구는 전체 확률 분포 $P(x)$ 가 아닌 국소 에러 $L_{X,Z}$ 를 손실 함수로 사용함으로써, 이 논문이 지적한 훈련 불가능성(집중 현상)을 완벽히 우회합니다.
    \end{itemize}

    \item \textbf{Architectures for Quantum Simulation Showing a Quantum Speedup}
    \begin{itemize}[label=--]
        \item \textbf{내용:} 얕은 깊이의 2D 격자 구조만으로도 충분한 얽힘을 생성하며, 단순한 비적응형 측정(Non-adaptive measurement)을 통해 고전 컴퓨터를 압도하는 양자 우위를 달성할 수 있음을 수학적으로 증명했습니다.
        \item \textbf{적용:} IBM Nighthawk와 같은 2D 사각 격자(Square Lattice) 프로세서의 물리적 토폴로지에 맞춘 SWAP-free 안자츠 설계의 근간이 됩니다.
    \end{itemize}
\end{itemize}

\subsection*{2. 본 연구의 방향성 및 차별점 요약}

본 연구의 핵심은 기존 연구들이 겪고 있는 \textbf{'이론적 훈련 불가능성'과 '하드웨어의 회로 깊이 폭발'이라는 딜레마를 훈련과 배포의 분리(TCDQ 철학) 및 국소 에러 기반 압축(BP-PPS)으로 동시에 타파}하는 데 있습니다.

\begin{itemize}
    \item \textbf{수학적 한계(Concentration)의 영리한 우회:} 타겟 데이터의 전체 확률 분포를 직접 학습하려다 바렌 플래토에 빠지는 기존 생성 모델들의 실패를 답습하지 않습니다. 고전 컴퓨터에서는 국소 관측가능량의 차이만 추적하여 파라미터 $\theta^*$ 를 효율적으로 훈련하고, 실제 복잡한 전역적 얽힘(Global entanglement) 샘플링은 양자 하드웨어에 온전히 맡깁니다.
    \item \textbf{IQP를 넘어선 범용 양자 동역학 모델링:} 제한적인 IQP 모델을 넘어서 2D 스핀 유리 모델의 비가환적 해밀토니안을 따르는 장시간(Long-time) 동역학을 타겟으로 삼습니다. 수백 스텝이 필요한 트로터 회로를 얕은 깊이($50$ 이하)의 HEA로 압축함으로써 하드웨어 결어긋남(Decoherence) 한계를 극복합니다.
    \item \textbf{다차원적 성능 비교 및 결과 검증 (차별점):} 단순히 얕은 회로를 실행하는 것에 그치지 않고, SOTA 고전 텐서 네트워크(PEPS 등) 및 일반적인 Trotter 방식과의 압축 효율성 비교를 수행합니다. 또한, 하드웨어에서 샘플링된 비트스트링 데이터가 타겟 분포의 물리적 특성을 띠는지 증명하는 교차 검증 알고리즘을 통해, 실용적이고 신뢰할 수 있는 양자 우위(Quantum Advantage)를 선언하는 것이 본 연구의 궁극적인 차별성입니다.
\end{itemize}

\vspace{0.6em}
\section*{II. 연구의 핵심 방향성 (Research Direction)}
\addcontentsline{toc}{section}{II. 연구의 핵심 방향성 (Research Direction)}

본 연구의 궁극적인 목표는 2D 스핀 유리(Spin Glass) 모델의 장시간 동역학(Dynamics)과 바닥 상태(Ground State)를 얕은 깊이의 양자 회로로 압축하여 샘플링하는 것입니다. 다만 두 과제의 성격은 다릅니다. 고전적으로 시뮬레이션이 불가능한 쪽은 장시간 실시간 동역학과 그 출력 분포 샘플링이며, 바닥 상태 준비는 (아래 1절에서 보이듯) 고전적으로 접근 가능하므로 회로 깊이 우위를 보이는 과제로 설정합니다. 이를 위해 먼저 타겟 물리계인 2D 스핀 유리 모델의 수학적 구조를 명확히 정의하고, 이를 바탕으로 세부적인 연구 수행 전략을 제시합니다.

\subsection*{1. 타겟 물리계: 2D 스핀 유리 (Spin Glass) 모델의 수학적 이해}

2D 스핀 유리 모델은 스핀 간의 상호작용이 무작위로 섞여 있어, 고전적인 최적화 및 샘플링 알고리즘이 임계 느려짐(Critical slowing down)에 빠지는 대표적인 강상관 양자 다체계입니다. 이 시스템의 해밀토니안 $H$ 는 다음과 같이 정의됩니다.

\begin{equation*}
H = -\sum_{\langle i,j \rangle} J_{ij} Z_i Z_j - h \sum_i X_i
\end{equation*}

각 수식 기호의 물리적 의미는 다음과 같습니다.

\begin{itemize}
    \item $\bm{\langle i,j \rangle}$ \textbf{:} 2D 사각 격자 상에서 서로 인접한(Nearest-neighbor) 큐비트 쌍을 의미합니다.
    \item $\bm{Z_i}$ \textbf{와} $\bm{X_i}$ \textbf{:} 각각 $i$ 번째 큐비트에 작용하는 파울리 연산자입니다. 파울리 $Z$ 는 스핀의 방향(위/아래)을 결정하고, 파울리 $X$ 는 스핀을 뒤집는 양자 요동(Quantum fluctuation)을 발생시킵니다.
    \item $\bm{h}$ \textbf{:} 모든 스핀에 균일하게 가해지는 횡방향 자기장(Transverse field)의 세기입니다.
    \item $\bm{J_{ij}}$ \textbf{(핵심 변수):} 인접한 스핀 간의 결합 상수입니다. 일반적인 강자성(Ferromagnetic) 모델에서는 모든 $J_{ij}$ 가 양수이지만, \textbf{스핀 유리 모델에서는 $J_{ij}$ 가 무작위적인 부호(예: $+1$ 또는 $-1$ )와 크기를 가집니다.}
\end{itemize}

\vspace{0.2em}
\textbf{고전적 샘플링이 불가능한 이유 (Frustration과 얽힘):} \\
스핀 유리 모델의 핵심은 $J_{ij}$ 의 무작위성으로 인해 발생하는 \textbf{쩔쩔맴(Frustration)} 현상에 있습니다. 예를 들어, 격자 내의 닫힌 루프(Loop)에서 결합 상수 $J_{ij}$ 의 곱이 음수가 되면, 시스템 내의 모든 스핀이 동시에 상호작용 에너지를 최소화하는 방향으로 정렬하는 것이 수학적으로 불가능해집니다.

이러한 Frustration과 횡방향 자기장 $h$ 의 양자 요동이 결합하면, 에너지가 비슷한 수많은 바닥 상태 후보군(Highly degenerate landscape)이 형성되고 스핀들 사이에 복잡한 장거리 얽힘(Long-range entanglement)이 발생합니다.

Frustration과 부호 문제(Sign problem)의 구분:
여기서 한 가지를 정확히 해 둘 필요가 있습니다. 계산 기저에서 $J_{ij} Z_i Z_j$ 항은 전부 대각 성분이며, 비대각 성분은 오직 $-h X_i$ 에서만 나오고 그 값은 모두 $-h < 0$ 입니다. 따라서 이 해밀토니안은 $J_{ij}$ 의 부호를 어떻게 섞더라도 stoquastic 이며, $e^{-\beta H}$ 의 양자 몬테카를로(QMC) 가중치가 전부 양수여서 부호 문제가 발생하지 않습니다. Frustration은 대각 에너지 지형의 성질이고 부호 문제는 비대각 가중치의 성질이므로, 이 모델에서 둘은 별개입니다.

그렇다면 고전적 어려움은 어디에 있는가:
\begin{itemize}

\item 실시간 동역학 $e^{-iHt}$ (본 연구의 핵심 타겟): QMC는 허수시간 전용이며, 실시간에서는 가중치가 복소수가 되어 동역학적 부호 문제로 붕괴합니다. 텐서 네트워크(PEPS 등)는 시간에 따른 얽힘 증가로 본드 차원이 폭발하고, 파울리 전파(SPD)는 연산자 확산으로 파울리 문자열 수가 폭발합니다. 세 계열 모두 장시간에서 실패하며, 이것이 본 연구가 겨냥하는 지점입니다.
\item 얕은 2D 회로의 출력 분포 $P(x)$ 샘플링: Anticoncentration을 갖는 분포의 샘플링은 복잡도 이론적 근거(Bermejo-Vega et al., PRX 8, 021010)로 고전적으로 어렵습니다. 고전 컴퓨터는 국소 관측량은 계산해도 전체 비트스트링을 뽑아내지는 못합니다.
\item 바닥 상태 에너지: 위에서 본 대로 stoquastic이므로 QMC가 적용 가능합니다. Frustration으로 인한 임계 느려짐(Critical slowing down)이라는 휴리스틱 장벽은 있으나 복잡도 이론적 장벽은 아니며, parallel tempering 등으로 100 스핀 규모는 도달 가능합니다. 따라서 본 연구에서 바닥 상태 준비는 양자 우위가 아니라 회로 깊이 우위(동일 깊이 예산에서 Trotter화된 회로 대비 우수한 에너지)로 주장하며, 기준값은 QMC로 잡습니다.
\end{itemize}

\subsection*{2. 핵심 연구 수행 전략}

본 연구는 앞서 정의한 2D 스핀 유리 모델의 복잡성을 고전 컴퓨터의 훈련 효율성과 양자 하드웨어의 샘플링 능력을 결합하여 돌파합니다.

\subsubsection*{가. 시간 진화 회로 압축 (Time-Evolution Compression)}
깊이가 기하급수적으로 깊어지는 트로터(Trotter) 시간 진화 연산자 $V = e^{-iHt}$ 를 깊이 $50$ 이하의 하드웨어 효율적 안자츠(HEA) $U(\theta)$ 로 압축합니다.

\begin{itemize}
    \item \textbf{손실 함수 정의:} 전역 확률 분포가 아닌, 국소적 파울리 연산자의 진화 오차를 측정하는 손실 함수 $\mathcal{L}_{X,Z}$ 를 사용합니다. 이는 목표 유니터리 $V$ 와 변분 안자츠 $U(\theta)$ 를 통과한 개별 큐비트의 파울리 $X$ 와 $Z$ 에러의 합으로 정의됩니다.
    \begin{equation*}
    \mathcal{L}_{X,Z} = \sum_i \|X_i(t) - \tilde{X}_i(t)\|_{\mathrm{rHS}}^2 + \sum_i \|Z_i(t) - \tilde{Z}_i(t)\|_{\mathrm{rHS}}^2
    \end{equation*}
    \item \textbf{BP-PPS 최적화:} 위 수식과 같이 국소 에러를 손실 함수로 사용할 경우, BP-PPS(Backpropagating Pauli Propagation) 알고리즘을 통해 파울리 문자열의 절사(Truncation) 오차를 통제하며 고전 컴퓨터에서 $\theta^*$ 에 대한 그래디언트를 메모리 효율적으로 빠르게 계산할 수 있습니다.
\end{itemize}

\subsubsection*{나. 바닥 상태 준비 및 훈련 (Ground State Preparation)}
동역학 시뮬레이션뿐만 아니라, 시스템의 바닥 상태를 찾는 문제에도 동일한 방법론을 적용합니다.

\begin{itemize}
    \item \textbf{에너지 최소화:} 해밀토니안 $H$ 의 기댓값인 에너지 $E = \langle H \rangle$ 를 최소화하는 방향으로 BP-PPS 기반의 고전적 경사하강법(Gradient Descent)을 수행합니다.
    \begin{equation*}
    E(\theta) = -\sum_{\langle i,j \rangle} J_{ij} \langle Z_i Z_j \rangle_{\theta} - h \sum_i \langle X_i \rangle_{\theta}
    \end{equation*}
    \item 이 과정 역시 국소 관측가능량인 $\langle Z_i Z_j \rangle$ 와 $\langle X_i \rangle$ 만을 평가하므로, 바렌 플래토(Barren Plateaus)에 빠지지 않고 안정적인 고전 훈련이 가능합니다.
\end{itemize}

\subsubsection*{다. 하드웨어 배포 및 양자 스크램블링 샘플링 (Deploy on Quantum)}
고전 컴퓨터는 파라미터 $\theta^*$ 를 효율적으로 찾아낼 수 있지만, 완성된 안자츠 $U(\theta^*)$ 가 만들어내는 전체 스핀 확률 분포 $P(x)$ 를 샘플링하는 것은 불가능합니다.

\begin{itemize}
    \item \textbf{양자 하드웨어 실행:} 찾아낸 최적 파라미터 $\theta^*$ 를 IBM 2D 사각 격자 하드웨어에 주입하여 회로를 실행합니다.
    \item \textbf{양자 우위 달성:} 깊이가 얕은 회로를 사용하므로 하드웨어 노이즈에 의한 결어긋남(Decoherence) 이전에 연산이 완료됩니다. 이를 통해 고전 시뮬레이터로는 닿을 수 없는 장시간 영역에서의 양자 정보 스크램블링 현상(예: OTOC(Out-of-Time-Order Correlator) 붕괴)을 관측하고, 시간진화된 상태 $U(\theta^*)|s\rangle$ 의 비트스트링 분포를 샘플링하여 실질적인 양자 생성 이득(Generative Quantum Advantage)을 입증합니다.
    \item \textbf{바닥 상태 샘플링은 우위 주장에서 제외:} Stoquastic 해밀토니안의 바닥 상태는 Perron--Frobenius 정리에 의해 계산 기저에서 모든 진폭을 음이 아니게 잡을 수 있고 (수치 확인: $3\times3$ 격자에서 최소 진폭 $+4.0\times10^{-3}$, 음수 진폭 0개), 그러면 $P(x)=|\psi_0(x)|^2$ 는 사실상 고전 확률분포여서 QMC가 직접 샘플링할 수 있습니다. 반면 시간진화된 상태는 복소 진폭을 가지며(같은 격자에서 위상 분포 폭 $\approx 2\pi$), 이것이 고전 샘플링을 막는 실제 원인입니다.
\end{itemize}

\vspace{0.6em}
\section*{III. 다른 연구와의 차별점 (Novelty \& Differentiation)}
\addcontentsline{toc}{section}{III. 다른 연구와의 차별점 (Novelty \& Differentiation)}

본 연구는 기존 양자 생성 모델과 양자 시뮬레이션 연구들이 지닌 이론적, 구조적 한계를 영리하게 우회합니다. 특히 IQP 회로의 제약을 벗어나 범용 동역학을 다루면서도, 훈련 불가능성(Barren Plateaus)을 국소적 최적화로 타파하여 확고한 양자 이득을 확보한다는 점에서 뚜렷한 차별성을 가집니다.

\subsection*{1. IQP의 구조적 한계 탈피 및 비가환적(Non-commuting) 범용 동역학 시뮬레이션}
\begin{itemize}
    \item 기존의 TCDQ(Train on Classical, Deploy on Quantum) 모델은 고전적 훈련의 효율성을 극대화하기 위해, 서로 교환 가능한(Commuting) 대각 게이트로만 구성된 IQP 회로를 사용했습니다.
    \item 그러나 이러한 IQP 구조는 2D 스핀 유리 모델처럼 국소 자기장과 스핀 상호작용이 뒤섞인 비가환적인 물리계의 역학을 온전히 표현할 수 없습니다.
    \item 본 연구는 IQP를 과감히 배제하고, IBM Nighthawk와 같은 2D 사각 격자 토폴로지에 최적화된 하드웨어 효율적 안자츠(HEA)를 도입합니다.
    \item 대신 BP-PPS(Backpropagating Pauli Propagation) 알고리즘을 적용하여, 비가환적 범용 회로라 하더라도 중간 상태 저장 없이 $O(n_{\mathrm{param}})$ 의 메모리 효율성으로 정확한 그래디언트를 계산하여 얕은 회로의 훈련을 가능하게 합니다.
\end{itemize}

\subsection*{2. 국소 에러 기반 학습을 통한 Barren Plateaus의 근본적 회피}
\begin{itemize}
    \item 생성 모델 연구에서 밝혀졌듯, 복잡한 난수 상태(Anticoncentration)를 다룰 때 전체 확률 분포를 타겟으로 삼는 MMD 등의 전역적(Global) 손실 함수를 사용하면 평균과 분산이 기하급수적으로 소실되는 바렌 플래토 현상이 발생하여 훈련 자체가 불가능해집니다.
    \item 본 연구는 이러한 전체 분포 학습의 한계를 인정하고, 대신 진화 연산에 따른 단일 큐비트 관측가능량(예: $X$, $Z$ )의 차이인 국소 에러 $L_{X,Z}$ 를 손실 함수로 설정합니다.
    \item 파울리 전파를 활용하여 국소적인 에러만을 역전파하는 이 방식은 깊이가 얕은 안자츠에서 기울기의 분산 소실을 수학적으로 방지하므로, 바렌 플래토를 완벽히 우회하여 고전 컴퓨터에서의 성공적인 최적화를 보장합니다.
    \item 파라미터 초기화 등 TCDQ 방법론 적용 고려
\end{itemize}

\subsection*{3. 훈련과 샘플링의 정보 비대칭성을 활용한 실질적 양자 이득(Quantum Advantage) 획득}
\begin{itemize}
    \item 고전 컴퓨터는 100 큐비트 급 시스템의 전체 얽힘을 추적할 수 없으므로, 국소 정보만을 절사(Truncation)하여 파라미터 최적화를 수행합니다.
    \item 그러나 훈련이 끝난 파라미터가 적용된 양자 하드웨어는 깊이 50 이하의 얕은 회로 내에서도 전역적 얽힘(Global entanglement) 상태를 자연스럽게 형성하고 유지합니다.
    \item 결과적으로 텐서 네트워크(PEPS 등)와 같은 고전 알고리즘으로는 본드 차원 폭발로 인해 시뮬레이션할 수 없는 장시간 스크램블링(OTOC) 현상이나 복잡한 2D 스핀 유리의 확률 분포를 양자 하드웨어의 샘플링만으로 성공적으로 추출해 냅니다. 이는 고전적 훈련의 효율성과 양자적 실행 복잡성을 분리하여 확정적인 양자 이득을 입증하는 본 연구만의 독창적 성과입니다.
\end{itemize}

\vspace{0.6em}
\section*{IV. 세부 연구 수행 단계 (Research Phases)}
\addcontentsline{toc}{section}{IV. 세부 연구 수행 단계 (Research Phases)}

본 연구는 고전 컴퓨터의 훈련 효율성과 양자 컴퓨터의 얽힘 샘플링 능력을 극대화하기 위해 총 5개의 페이즈(Phase)로 나누어 진행됩니다.

\subsection*{Phase 1: 안자츠(Ansatz) 구축 및 비교 모델(Trotter, 고전) 준비}
첫 번째 단계는 IBM Nighthawk 프로세서의 물리적 토폴로지를 완벽히 반영한 양자 회로와 벤치마킹을 위한 대조군을 설계하는 것입니다.

\begin{itemize}
    \item \textbf{HEA (Hardware-Efficient Ansatz) 구축:} 사각 격자(Square Lattice) 구조에 맞추어, 인접한 큐비트 간의 상호작용만 허용하여 SWAP 게이트 오버헤드를 완전히 제거한 깊이 $50$ 이하의 얕은 변분 회로를 설계합니다.
    \item \textbf{손실 함수(Loss Function) 구체화:} BP-PPS 논문에 기반하여, 목표로 하는 동역학 연산자의 파울리 전파 에러를 최소화하는 국소 에러 손실 함수
    \begin{equation*}
    \mathcal{L}_{X,Z} = \sum_{i} \|X_i(t) - \tilde{X}_i(t)\|_{\mathrm{rHS}}^2 + \sum_{i} \|Z_i(t) - \tilde{Z}_i(t)\|_{\mathrm{rHS}}^2
    \end{equation*}
    를 정의합니다.
    \item \textbf{비교 모델 준비:} 양자 하드웨어의 압축 이점을 비교하기 위해, 타겟 동역학에 대한 다항식 차수의 트로터화(Trotterization) 회로를 준비합니다. 아울러 고전 시뮬레이션의 한계를 보여줄 텐서 네트워크(PEPS 등) 모델을 세팅합니다.
\end{itemize}

\subsection*{Phase 2: 고전적 파라미터 최적화(Training) 방법론 구체화}
BP-PPS(Backpropagating Pauli Propagation) 알고리즘을 활용하여 고전 컴퓨터 내에서 파라미터 $\theta$ 를 효율적으로 훈련합니다.

\begin{itemize}
    \item \textbf{최적화 알고리즘 적용:} BP-PPS는 정확한 해석적 그래디언트를 제공하므로, 2차 미분 정보를 근사하여 빠르고 정밀하게 수렴하는 \textbf{L-BFGS-B} 옵티마이저를 기본으로 채택합니다. 또한, 초기 탐색 공간이 넓고 지역 최솟값(Local minima)에 빠질 위험을 줄이기 위해, 모멘텀(Momentum) 기법이 결합된 \textbf{Adam} 옵티마이저를 훈련 초기에 활용하는 하이브리드 미세 조정(Fine-tuning) 전략을 도입합니다.
    \item \textbf{절사(Truncation) 오차 관리:} 고전 메모리 폭발을 막기 위해 파울리 문자열의 계수 임계값 $\delta$ 를 설정하여 의미 없는 항을 버립니다. 최적화 진행 상황에 따라 임계값 $\delta$ 를 동적으로 조절하여 훈련 정확도와 메모리 사용량의 트레이드오프를 엄밀하게 관리합니다.
    \item Pauli Weight 등의 다른 방법론 적용 여부도 고려 가능합니다.
\end{itemize}

\subsection*{Phase 3: 훈련 성능 비교와 검증 (Benchmarking on Classical)}
양자 하드웨어에 배포하기 전, 고전 컴퓨터 상에서 압축된 회로의 훈련 성능을 타 모델들과 철저히 벤치마킹합니다.

\begin{itemize}
    \item \textbf{Trotter 모델 대비 압축률 검증:} 훈련된 HEA 모델이 목표 시간 $t$ 에 도달하는 데 필요한 회로 깊이를 2차 및 4차 Trotter 회로의 깊이와 비교합니다. 동등한 국소 관측가능량 오차를 기준으로 회로 깊이가 얼마나 획기적으로(예: 40\% 이상) 감소했는지 정량적으로 분석합니다.
    \item \textbf{PEPS 최적화 속도 비교:} 고전 텐서 네트워크 최적화 방식 대비 BP-PPS 기반의 파라미터 수렴 속도 및 계산 복잡도(Runtime) 우위를 검증합니다.
\end{itemize}

\subsection*{Phase 4: 양자 하드웨어 배포 및 실행 (Deploy on Quantum)}
훈련이 완료된 최적의 파라미터 $\theta^*$ 를 기반으로 실제 양자 하드웨어에서 샘플링을 수행합니다.

\begin{itemize}
    \item \textbf{IBM Nighthawk 배포:} 최적화된 얕은 회로를 100 큐비트 규모의 2D 사각 격자 양자 프로세서에 배포합니다.
    \item \textbf{동역학 및 스크램블링 실행:} 회로가 하드웨어 결어긋남(Decoherence) 한계에 도달하기 전인 상수 시간 내에 장시간 스크램블링의 척도가 되는 OTOC(Out-of-Time-Order Correlator) 회로를 실행하여 복잡한 양자 상태의 얽힘 데이터를 생성합니다.
\end{itemize}

\subsection*{Phase 5: 샘플링 데이터 검증 알고리즘 및 양자 우위 입증 (Verification \& Validation)}
양자 하드웨어에서 얻은 데이터가 단순한 노이즈가 아님을 입증하고 최종적인 양자 우위를 선언합니다.

\begin{itemize}
    \item \textbf{비트스트링 검증 알고리즘:} 양자 장치에서 생성된 샘플들이 물리적으로 타당한지 확인하기 위해, 고전 그림자(Classical Shadows) 기법이나 특정 국소 에너지 기댓값 측정을 통해 고전적으로 미리 계산된 타겟 값과 오차 범위 내에서 일치하는지 교차 검증(Cross-validation)합니다.
    \item \textbf{샘플링 정확도 비교:} 특정 시간 $t$ 가 지난 후, 고전 텐서 네트워크(PEPS)가 본드 차원(Bond Dimension) 폭발로 인해 정확한 분포를 샘플링하지 못하고 무너지는 지점을 찾습니다.
    \item \textbf{양자 우위 선언:} 고전 시뮬레이터가 실패하는 바로 그 지점에서, BP-PPS로 얕게 압축된 우리의 양자 회로가 일관성 있는 OTOC 감쇠와 고품질의 비트스트링을 뱉어냄을 보여, 실질적이고 경험적인 양자 생성 우위(Empirical Generative Quantum Advantage)를 최종적으로 입증합니다.
\end{itemize}

\end{document}
